% AI review of the deck: reviewer setup, review prompt, resulting grade.
% Grade macros (\resReviewScore, \resReviewVerbal) come from review.tex,
% maintained by the orchestrator.
% Grade of the AI presentation review, copied verbatim by the orchestrator
% from the presentation-reviewer output (final iteration, 2026-10-20).
% Iterations are documented in the AI work log. This file is data for the
% grade slide; do not edit by hand on the slide side.
\newcommand{\resReviewScore}{9}
\newcommand{\resReviewVerbal}{Prezentacja spełnia wszystkie warunki
obowiązkowe, ma zwięzły wstęp z poprawnymi wzorami EMA, MACD i SIGNAL oraz
rzetelnie udokumentowaną współpracę z AI. Tytuły slajdów są tezami popartymi
treścią, a układ jest czytelny i spójny na pełnym ekranie. Wstęp ma
wyłącznie tekstowy charakter, więc przykładowy wykres przecięć MACD i
SIGNAL ułatwiłby jego zrozumienie.}


\begin{frame}
  \frametitle{Prezentację ocenia osobny agent z trzema źródłami kryteriów}
  \begin{itemize}
    \item Agent presentation-reviewer czyta skompilowany PDF i źródła LaTeX;
      ma wyłącznie narzędzia odczytu.
    \item Źródła kryteriów: zasady zaliczenia, instrukcja projektu, przewodnik
      stylu prezentacji (tytuły w formie tez).
    \item Ocenia: elementy obowiązkowe (brak to blocker), zrozumiałość
      wstępu, czytelność na pełnym ekranie, poprawność wzorów i wykresów,
      zgodność wniosków z dowodami, objętość względem 8 minut.
    \item Recenzja jest iteracyjna: każda iteracja daje znaleziska z wagami
      oraz ocenę słowną i liczbową, trafia do dziennika pracy i jest
      doradcza, nie przesądza o punktach.
  \end{itemize}
\end{frame}

\begin{frame}[fragile]
  \frametitle{Recenzję uruchamia prompt wskazujący źródła kryteriów i skalę ocen}
  {\footnotesize
\begin{verbatim}
Review slides/preliminary/main.pdf and its .tex sources
against: (1) [...]zasady_zaliczenia_projektu_v1.pdf, mandatory
conditions and slide groups; (2) [...]Projekt_1.pdf
table 1, preliminary deck content [...];
(3) docs/presentation-style.md, assertion-evidence titles,
budgets, palette, banned elements. Report findings with
severity and a suggested fix. End with a verbal grade
and a score 1-10.
\end{verbatim}}
  \vskip1mm
  Wynik: lista znalezisk z wagami i poprawkami oraz ocena słowna i liczbowa.\\
  Ocena: wskazanie trzech źródeł kryteriów zawęziło recenzję do wymagań kursu.
\end{frame}

\begin{frame}
  \frametitle{Recenzent AI ocenił prezentację wstępną na \resReviewScore{}/10}
  „\resReviewVerbal''
  \vskip3mm
  {\footnotesize\color{mnMuted}
    Ocena pochodzi z końcowej iteracji recenzji; przebieg iteracji
    dokumentuje dziennik pracy z AI.\par}
\end{frame}
